
          %%%%% ~~~~~~~~~~~~~~~~~~~~ %%%%%

\chapter{Blackbody radiation}
\setcounter{theorem}{0}
\setcounter{equation}{0}

%\cite{vanDerVaart1996}
%\cite{Kosorok2008}

%\renewcommand{\theenumi}{\alph{enumi}}
%\renewcommand{\labelenumi}{\textnormal{(\theenumi)}$\;\;$}
\renewcommand{\theenumi}{\roman{enumi}}
\renewcommand{\labelenumi}{\textnormal{(\theenumi)}$\;\;$}

          %%%%% ~~~~~~~~~~~~~~~~~~~~ %%%%%

\section{Universality of the energy density spectrum of blackbody radiation -- Kirchhoff's law of thermal radiation}

          %%%%% ~~~~~~~~~~~~~~~~~~~~ %%%%%

\vskip 0.5cm
\begin{definition}[Blackbody radiation cavity]
\mbox{}
\vskip 0.1cm
\noindent
A \textbf{blackbody radiation cavity} is, by definition,
a finite \textbf{\color{red}evacuated} cavity
with perfectly reflecting walls, maintained
at a uniform absolute temperature \,$T$,\,
where the radiation inside is in thermal equilibrium with the walls.
We remark that blackbody radiation cavities are
characterized by their respective macroscopic configurations
(encompassing variations in shape, volume, and interior wall material, etc.).
\end{definition}

\vskip 0.2cm
\begin{theorem}[Universality of the energy density spectrum of blackbody radiation]
\mbox{}
\vskip 0.1cm
\noindent
Suppose:
\begin{itemize}
\item
    $\mathcal{C}$\, is the set of all possible blackbody radiation cavities.
\item
    $T \in \Re^{+}$ is the absolute temperature, and
    \,$\nu \in \Re^{+}$ is the frequency of electromagnetic radiation.
\item
    \,$\rho : \Re^{+} \times \Re^{+} \times \mathcal{C} \xrightarrow{{\color{white}....}} \Re^{+}$\,
    is the spectral energy density function,
    i.e., \,$\rho(\nu, T, C)$\, is the energy density per unit frequency
    (energy per unit volume per unit frequency)
    at frequency \,$\nu \in \Re^{+}$,\, temperature \,$T \in \Re^{+}$,\,
    for the blackbody radiation cavity \,$C \in \mathcal{C}$.\,
\end{itemize}
Assume:
\begin{itemize}
\item
    \textbf{Axiom 1 (Continuity):}\;
    For any fixed $T$ and $C$, the spectral energy density
    \,$\rho(\nu, T, C)$\,
    is a continuous function of the frequency $\nu$.
\item
    \textbf{Axiom 2 (Strict Monotonicity of Flux):}\;
    The spectral power per unit area \,$\Phi_\nu$\,
    (energy per unit time per unit area)
    emitted from a small aperture in cavity \,$C$\,
    is a strictly monotonically increasing function
    of its internal spectral energy density.
    That is,
    \,$\Phi_{\nu} \, = \, f(\rho(\nu, T, C))$,\,
    where
    \,$f: \Re^{+} \to \Re^{+}$\,
    satisfies
    \,$f(a) \, > \, f(b) \;\Longleftrightarrow\; a \, > \, b$.\,
    (In classical electrodynamics,
    \,$f(u) \,=\, \dfrac{c}{4}u$,\,
    but strict monotonicity is mathematically sufficient).
\item
    \textbf{Axiom 3 (Second Law of Thermodynamics):}\;
    If two blackbody radiation cavities
    \,$C_{1}, C_{2} \in \mathcal{C}$\,
    are maintained at the exact same temperature \,$T$\, and
    are allowed to exchange energy radiatively and only radiatively,
    then the net steady-state heat flux between them must be exactly zero.
\item
    \textbf{Axiom 4 (Ideal Optical Filtration):}\;
    For any interval
    \,$[\,\nu_{1}, \nu_{2}\,] \subset \Re^{+}$,\,
    it is possible to construct a filter that transmits \,$100\%$\,
    of the incident radiation with frequencies lying within the interval
    \,$[\,\nu_{1}, \nu_{2}\,]$\,
    and reflects \,$100\%$\, of the incident radiation
    with frequencies lying outside this interval.
\end{itemize}
Then, for any two blackbody radiation cavities
\,$C_{1}, C_{2} \,\in\, \mathcal{C}$\,
(with arbitrary -- possibly distinct -- respective macroscopic configurations),
their respective spectral energy densities are equal
at each frequency and at each temperature, i.e.,
\begin{equation*}
\rho(\nu, T, C_{1})
\;\;=\;\;
    \rho(\nu, T, C_{2}),
\quad
\textnormal{for each \,$\nu > 0$,\, each \,$T > 0$,\, and each \,$C_{1}, C_{2} \in \mathcal{C}$}
\end{equation*}
Consequently,
\,$\rho(\,\nu,T,C\,)$\,
is in fact independent of \,$C$\, and
induces a universal function \,$\rho(\,\nu,T\,)$.\,
\end{theorem}
\proof
Suppose, to the contrary, that there exist a specific temperature
\,$T_{0} > 0$,\, a frequency \,$\nu_0 > 0$,\, and two cavities
\,$C_{1}, C_{2} \in \mathcal{C}$\,
such that their spectral energy densities differ
at frequency \,$\nu_{0}$\, and temperature \,$T_{0})$.\,

\vskip 0.2cm
\noindent
Without loss of generality, we may assume:
\begin{equation*}
\rho(\nu_0, T_{0}, C_{1})
\;\; > \;\;
    \rho(\nu_0, T_{0}, C_{2})
\end{equation*}
Because \,$u$\, is continuous with respect to \,$\nu$\, (Axiom 1),
there must exist a finite frequency neighbourhood \,$[\,\nu_{1}, \nu_{2}\,]$\,
containing \,$\nu_{0}$\, (\,where \,$\nu_{1} < \nu_{2}$\,)
such that the following inequality holds:
\begin{equation*}
\rho(\nu, T_{0}, C_{1})
\;\; > \;\;
    \rho(\nu, T_{0}, C_{2})\,,
\quad
\textnormal{for each \,$\nu \in [\,\nu_{1}, \nu_{2}\,]$}
\end{equation*}

\vskip 0.2cm
\noindent
Let the two cavities
\,$C_{1}$\, and \,$C_{2}$\,
be placed facing each other,
connected only by a small aperture of area \,$A > 0$.\,
Between the cavities, insert an ideal filter (Axiom 4)
that exclusively transmits radiation in the frequency range \,$[\,\nu_{1}, \nu_{2}\,]$,\,
reflecting all other frequencies back into their respective cavities.
The total power \,$P_{1 \to 2}$\, radiated from cavity \,$C_{1}$\, to \,$C_{2}$\,
through the filter is the integral of the emitted energy flux over the transmitted frequencies:
\begin{equation*}
P_{1 \rightarrow 2}
\;\; = \;\;
    A \int_{\nu_{1}}^{\nu_{2}} f(\rho(\nu, T_{0}, C_{1})) \, \d\nu
\end{equation*}

\vskip 0.2cm
\noindent
Similarly, the total power radiated from cavity \,$C_{2}$\, to \,$C_{1}$\, is:
\begin{equation*}
P_{2 \rightarrow 1}
\;\; = \;\;
    A \int_{\nu_{1}}^{\nu_{2}} f(\rho(\nu, T_{0}, C_{2})) \, \d\nu
\end{equation*}

\vskip 0.2cm
\noindent
By Axiom 2, since
\,$\rho(\nu, T_{0}, C_{1}) \,>\, \rho(\nu, T_{0}, C_{2})$\,
for each 
\,$\nu \in [\nu_{1}, \nu_{2}]$,\,
it follows that
\,$f(\rho(\nu, T_{0}, C_{1})) \,>\, f(\rho(\nu, T_{0}, C_{2}))$\,
for each 
\,$\nu \in [\,\nu_{1}, \nu_{2}\,]$,\,
which in turn implies:
\begin{equation*}
P_{1 \rightarrow 2}
\;\; = \;\;
    A \int_{\nu_{1}}^{\nu_{2}} f(\rho(\nu, T_{0}, C_{1})) \, \d\nu
\;\; > \;\;
    A \int_{\nu_{1}}^{\nu_{2}} f(\rho(\nu, T_{0}, C_{2})) \, \d\nu
\;\; = \;\;
    P_{2 \rightarrow 1}
\end{equation*}
Thus, there is a net transfer per unit time of heat energy
\,$P_{\textnormal{net}} \,=\, P_{1 \to 2} - P_{2 \to 1} \,>\, 0$\,
flowing continuously from cavity $C_{1}$ to cavity $C_{2}$.

\vskip 0.2cm
\noindent
However, both \,$C_{1}$\, and \,$C_{2}$\, are
at precisely the same temperature \,$T_{0}$.\,
A spontaneous, continuous net flow of heat between two bodies
at identical temperatures constitutes a violation of
the Second Law of Thermodynamics (Axiom 3).

\vskip 0.2cm
\noindent
Therefore, our initial contrary assumption must be false.
We may now conclude that no such
\,$\nu_0 > 0$\, and \,$T_{0} > 0$\,exist,
and hence,
\begin{equation*}
\rho(\nu, T, C_{1})
\;\; = \;\;
    \rho(\nu, T, C_{2})\,,
\quad
\textnormal{for each \,$\nu,\, T \,>\, 0$\, and each \,$C_{1}, C_{2} \,\in\, \mathcal{C}$}
\end{equation*}
This completes the proof of the Theorem.
\qed

          %%%%% ~~~~~~~~~~~~~~~~~~~~ %%%%%

\clearpage
\section{Blackbody radiation spectral energy density}

          %%%%% ~~~~~~~~~~~~~~~~~~~~ %%%%%

\vskip 0.3cm
\begin{proposition} %[Motivation for \,$\rho(\nu,T)\; = \; g(\nu) \cdot \langle\,E(\nu,T)\,\rangle$]
\mbox{}
\vskip 0.2cm
\noindent
Let
\begin{eqnarray*}
\rho(\,\nu,T\,)
{\color{white}.}
& := &
    \left(\begin{array}{c}
        \textnormal{universal}
        \\
        \textnormal{spectral energy density function}
        \\
        \textnormal{of blackbody radiation}
        \\
        \textnormal{at frequency \,$\nu$\, and temperature \,$T$}
        \end{array}\right),
\\
g(\nu)
{\color{white}....}
& := &
    \left(\begin{array}{c}
        \textnormal{number of all}
        \\
        \textnormal{admissible modes of oscillation}
        \\
        \textnormal{at frequency \,$\nu$,}
        \\
        \textnormal{per unit volume, per unit frequency}
        \end{array}\right),
    \quad
    \textnormal{and}
\\
\langle\,E(\nu,T)\,\rangle
& := &
    \left(\begin{array}{c}
        \textnormal{average energy over all}
        \\
        \textnormal{admissible modes of oscillation}
        \\
        \textnormal{at frequency \,$\nu$\, and temperature \,$T$}
        \end{array}\right)
\end{eqnarray*}
Then,
\begin{equation*}
{\color{red}
\rho(\nu,T)
\;\; = \;\;
    g(\nu) \langle\,E(\nu,T)\,\rangle}\,,
\quad
\textnormal{for each \,$\nu \in (\,0\,,\,\infty\,)$\, and each \,$T \in (\,0\,,\,\infty\,)$},
\end{equation*}
\end{proposition}
\proof
Let \,$\nu_{1}, \nu_{2} \in (\,0\,,\infty\,)$\,
with \,$\nu_{1} < \nu_{2}$,\, but otherwise arbitrary.
Note that:
% \begin{equation*}
% % N(\nu)
% % \;\; := \;\;
%     \left(\begin{array}{c}
%         \textnormal{number of all}
%         \\
%         \textnormal{admissible modes of oscillation}
%         \\
%         \textnormal{with frequencies between \,$\nu_{1}$\, and \,$\nu_{2}$}
%         \end{array}\right)
% \;\; = \;\;
%     \int_{\nu_{1}}^{\nu_{2}}\, V \cdot g(\nu) \,\d\nu
% \end{equation*}
% Next, observe that
\begin{equation*}
\left(\begin{array}{c}
    \textnormal{total amount of energy}
    \\
    \textnormal{at temperature \,$T$,}
    \\
    \textnormal{contained in}
    \\
    \textnormal{oscillators with frequencies}
    \\
    \textnormal{between \,$\nu_{1}$\, and \,$\nu_{2}$}
    \end{array}\right)
\;\; =: \;\;
    U_{\textnormal{total}}(\nu_{1},\nu_{2}\,;T)
% \;\; = \;\;
%     \int_{\nu_{1}}^{\nu_{2}} N(\nu) \langle\,E(\nu,T)\,\rangle \,\d\nu
\;\; = \;\;
    \int_{\nu_{1}}^{\nu_{2}} V \cdot g(\nu) \langle\,E(\nu,T)\,\rangle \,\d\nu
\end{equation*}
It follows that
\begin{equation*}
\int_{\nu_{1}}^{\nu_{2}}\, g(\nu) \langle\,E(\nu,T)\,\rangle \,\d\nu
\;\; = \;\;
    \dfrac{U_{\textnormal{total}}(\nu_{1},\nu_{2}\,;T)}{V}
\;\; = \;\;
    % \left(\begin{array}{c}
    %     \textnormal{energy}
    %     \\
    %     \textnormal{per unit volume}
    %     \end{array}\right)
    \left(\begin{array}{c}
        \textnormal{total amount of energy}
        \\
        \textnormal{per unit volume}
        \\
        \textnormal{at temperature \,$T$,}
        \\
        \textnormal{contained in}
        \\
        \textnormal{oscillators with frequenciess}
        \\
        \textnormal{between \,$\nu_{1}$\, and \,$\nu_{2}$}
        \end{array}\right)
\;\; = \;\;
    \int_{\nu_{1}}^{\nu_{2}}\,\rho(\nu,T)\,\d\nu
\end{equation*}
Hence,
\begin{equation*}
\int_{\nu_{1}}^{\nu_{2}}
    \left(\,
        \rho(\nu,T)
        \, \overset{{\color{white}1}}{-} \,
        g(\nu) \langle\,E(\nu,T)\,\rangle
        \right)
    \,\d\nu
\;\; = \;\;
    0
\end{equation*}
Lastly, recall that \,$\nu_{1}, \nu_{2} \in (\,0\,,\,\infty\,)$\,
satisfy
\,$\nu_{1} < \nu_{2}$,\,
but are otherwise arbitrary.
By continuity of the integrand, we may thus conclude that
\begin{equation*}
\rho(\nu,T)
\;\; = \;\;
    g(\nu) \langle\,E(\nu,T)\,\rangle\,,
\quad
\textnormal{for each \,$\nu \in (\,0\,,\,\infty\,)$\, and each \,$T \in (\,0\,,\,\infty\,)$},
\end{equation*}
as required.
\qed

\vskip 0.5cm
\begin{proposition} %[Motivation for \,$g(\nu) \,=\, 8\,\pi\,\nu^{2}\,/\,c^{3}$]
\mbox{}
\vskip 0.2cm
\noindent
Suppose:
\begin{itemize}
\item
    The Cavity Model (Thermal Equilibrium):
    The blackbody is modeled as an \textbf{\color{red}evacuated} cubic cavity
    of side length \,$L > 0$\, and hence volume \,$V = L^3$,\,
    with perfectly reflecting walls, maintained
    at a uniform absolute temperature \,$T$.\,
    The radiation inside is in thermal equilibrium with the walls.
\item
    Electromagnetic Standing Waves: The radiation field inside the cavity
    consists of independent electromagnetic standing waves.
    The electric field must vanish at the conducting boundaries of the cavity,
    constraining the allowed wavelengths (modes) of the system.
\end{itemize}
Let
\begin{equation*}
g(\nu)
\;\; := \;\;
    \left(\begin{array}{c}
        \textnormal{number of all}
        \\
        \textnormal{admissible modes of oscillation}
        \\
        \textnormal{per unit volume, per unit frequency,}
        \\
        \textnormal{at frequency \,$\nu$}
        \end{array}\right)
% \;\; = \;\;
%     \dfrac{8\,\pi}{c^{3}}\cdot\nu^{2}
\end{equation*}
Then,
\begin{equation*}
{\color{red}
g(\nu)
\;\; = \;\;
    \dfrac{
        8\,\pi
        }{
        c^{3}
        }
    \cdot
    \nu^{2}
}\,,
\quad
\textnormal{for each \,$\nu \in (\,0\,,\,\infty\,)$},
\end{equation*}
\end{proposition}
\proof

\vskip 0.3cm
\noindent
First, we define:
\begin{equation*}
N(\nu)
\;\; := \;\;
    \left(\begin{array}{c}
        \textnormal{number of all}
        \\
        \textnormal{admissible modes of oscillation}
        \\
        \textnormal{with frequencies {\color{red}less than or equal to \,$\nu$}}
        \end{array}\right)
% \;\; = \;\;
%     \int_{\nu_{1}}^{\nu_{2}}\, V \cdot g(\nu) \,\d\nu
\end{equation*}

\vskip 0.4cm
\noindent
\textbf{Claim 1:}\;
$
N(\nu)
\; = \;
    \dfrac{
        8\,\pi\,V\,\nu^{3}
        }{
        3\,c^{3}
        }
$
\vskip 0.1cm
\noindent
Proof of Claim 1:\;
Let
\,$\mathbf{k} \,=\, (\,k_{x}, k_{y}, k_{z}\,) \,\in\, \Re^{3}$\,
be the wave vector of any electromagnetic standing wave
inside the cavity.
The boundary conditions that the standing wave amplitude must
be zero at the walls of the cubic cavity imply:
\begin{equation*}
k_{x} \;=\; \dfrac{n_{x}\pi}{L},
\quad
k_{y} \;=\; \dfrac{n_{y}\pi}{L},
\quad
k_{z} \;=\; \dfrac{n_{z}\pi}{L},
\end{equation*}
where
\,$n_{x},\, n_{y},\, n_{z} \,\in\, \N$.\,
Hence,
\begin{eqnarray*}
\vert\,\mathbf{k}\,\vert^{2}
& = &
    \left(k_{x}\right)^{2}
    \,+\,
    \left(k_{y}\right)^{2}
    \,+\,
    \left(k_{z}\right)^{2}
\;\; = \;\;
    \left(\dfrac{n_{x}\pi}{L}\right)^{2}
    \,+\,
    \left(\dfrac{n_{y}\pi}{L}\right)^{2}
    \,+\,
    \left(\dfrac{n_{z}\pi}{L}\right)^{2}
\;\; = \;\;
    \left(\,\dfrac{\pi}{L}\,\right)^{2}
    \left(n_{x}^{2}+n_{y}^{2}+n_{z}^{2}\right)
\\
& = &
    \left(\,\dfrac{\pi}{L}\,\right)^{2}
    \vert\,\mathbf{n}\,\vert^{2}\,,
\end{eqnarray*}
where
\,$\mathbf{n} \,=\, (\,n_{x},n_{y},n_{z}\,)$.\,

\vskip 0.2cm
\noindent
Now, recall the relation between the magnitude of the wave vector \,$\mathbf{k}$\,
and the frequency \,$\nu$\, and the wavelength \,$\lambda \,=\, c/\nu$\,
of the (standing) wave:
\begin{equation*}
\vert\,\mathbf{k}\,\vert
\; = \;
    \dfrac{2\,\pi}{\lambda}
\; = \;
    \dfrac{2\,\pi}{c/\nu}
\; = \;
    \dfrac{2\,\pi\,\nu}{c}
\end{equation*}
Hence,
\begin{equation*}
\left(\,\dfrac{2\,\pi\,\nu}{c}\,\right)^{2}
\;\; = \;\;
    \vert\,\mathbf{k}\,\vert^{2}
\;\; = \;\;
    \left(\,\dfrac{\pi}{L}\,\right)^{2}
    \vert\,\mathbf{n}\,\vert^{2}\,,
\end{equation*}
which in turn implies:
\begin{equation*}
\vert\,\mathbf{n}\,\vert
\;\; = \;\;
    \dfrac{2\,L\,\nu}{c}
\end{equation*}
Now, recall the definition:
\begin{equation*}
N(\nu)
\;\; := \;\;
    \left(\begin{array}{c}
        \textnormal{number of all}
        \\
        \textnormal{admissible modes of oscillation}
        \\
        \textnormal{with frequencies {\color{red}less than or equal to \,$\nu$}}
        \end{array}\right)
\end{equation*}
Taking into account the fact that electromagnetic waves with the same wave vector
admit two polarizations, we see that
\,$N(\nu)$\,
is \,$2$\, times the number of lattice points with non-negative integer coordinates
contained inside the intersection of the first octant and
the sphere of radius \,$\vert\,\mathbf{n}\,\vert$.\,
Hence,
\begin{equation*}
N(\nu)
\;\; = \;\;
    2 \times \dfrac{1}{8} \times \left(\dfrac{4}{3}\,\pi\,\vert\,\mathbf{n}\,\vert^{3}\right)
\;\; = \;\;
    \dfrac{\pi}{3} \cdot \vert\,\mathbf{n}\,\vert^{3}
\;\; = \;\;
    \dfrac{\pi}{3} \cdot \left(\dfrac{2\,L\,\nu}{c}\right)^{3}
\;\; = \;\;
    \dfrac{8\,\pi\,L^{3}\,\nu^{3}}{3\,c^{3}}
\;\; = \;\;
    \dfrac{8\,\pi\,V\,\nu^{3}}{3\,c^{3}}
\end{equation*}
This completes the proof of Claim 1.

\vskip 0.4cm
\noindent
\textbf{Claim 2:}\;\,
$
% g(\nu)
\dfrac{1}{V}\cdot\dfrac{\d N(\nu)}{\d\,\nu}
\; = \;
    \dfrac{
        8\,\pi
        }{
        c^{3}
        }
    \cdot
    \nu^{2}
$
\vskip 0.2cm
\noindent
Proof of Claim 2:\;
By Claim 1, we have
\begin{equation*}
% g(\nu)
% \;\; = \;\;
\dfrac{1}{V}\cdot\dfrac{\d N(\nu)}{\d\,\nu}
\;\; = \;\;
    \dfrac{1}{V}
    \cdot
    \dfrac{\d}{\d\,\nu}\!\left(\,
        \dfrac{8\,\pi\,V\,\nu^{3}}{3\,c^{3}}
        \,\right)
\;\; = \;\;
    \dfrac{8\,\pi}{3\,c^{3}}
    \cdot
    \dfrac{\d}{\d\,\nu}\!\left(\,
        \nu^{3}
        \,\right)
\;\; = \;\;
    \dfrac{8\,\pi}{3\,c^{3}}
    \cdot
    3\,\nu^{2}
\;\; = \;\;
    \dfrac{8\,\pi}{c^{3}}
    \cdot
    \nu^{2}\,,
\end{equation*}
This completes the proof of Claim 2.

\vskip 0.4cm
\noindent
\textbf{Claim 3:}\;\,
$
g(\nu)
\; = \;
    \dfrac{1}{V}\cdot\dfrac{\d N(\nu)}{\d\nu}\,,
$\;\;
for each \,$\nu \in (\,0\,,\,\infty\,)$.
\vskip 0.2cm
\noindent
Proof of Claim 3:\;
Let \,$\nu_{1}, \nu_{2} \in (\,0\,,\infty\,)$\,
with \,$\nu_{1} < \nu_{2}$,\, but otherwise arbitrary.
Note that
\begin{equation*}
\int_{\nu_{1}}^{\nu_{2}}\, \dfrac{\d N(\nu)}{\d\nu} \,\d\nu
% \;\; = \;\;
%     N(\nu)
\;\; = \;\;
    \left(\begin{array}{c}
        \textnormal{number of all}
        \\
        \textnormal{admissible modes of oscillation}
        \\
        \textnormal{with frequencies between \,$\nu_{1}$\, and \,$\nu_{2}$}
        \end{array}\right)
\;\; = \;\;
    \int_{\nu_{1}}^{\nu_{2}}\, V \cdot g(\nu) \,\d\nu\,,
\end{equation*}
which implies:
\begin{equation*}
\int_{\nu_{1}}^{\nu_{2}}
\left(\;
    g(\nu)
    \; \overset{{\color{white}1}}{-} \;
    \dfrac{1}{V}\cdot\dfrac{\d N(\nu)}{\d\nu}
    \,\right)\,\d\nu
\;\; = \;\;
    0
\end{equation*}
Now, recall that \,$\nu_{1}, \nu_{2} \in (\,0\,,\,\infty\,)$\,
satisfy
\,$\nu_{1} < \nu_{2}$,\,
but are otherwise arbitrary.
By continuity of the integrand, we may thus conclude that
\begin{equation*}
g(\nu)
\;\; = \;\;
    \dfrac{1}{V}\cdot\dfrac{\d N(\nu)}{\d\nu}\,,
\quad
\textnormal{for each \,$\nu \in (\,0\,,\,\infty\,)$}
\end{equation*}
This completes the proof of Claim 3.

\vskip 0.4cm
\noindent
The Proposition now follows immediately from Claim 2 and Claim 3.
\qed

          %%%%% ~~~~~~~~~~~~~~~~~~~~ %%%%%

\vskip 0.5cm
\section{The ultraviolet catastrophe -- the Rayleigh–Jeans law}

          %%%%% ~~~~~~~~~~~~~~~~~~~~ %%%%%

\noindent
\textbf{The Rayleigh-Jeans assumptions}
\begin{itemize}
\item
    Inside a given blackbody radiation cavity,
    for each fixed frequency \,$\nu \in (\,0,\infty\,)$,\,
    the set \,$\mathscr{E}_{\nu}$\, of admissible energies
    among all modes of electromagnetic oscillation
    with frequency \,$\nu$\, is given by:
    \begin{equation*}
    \mathscr{E}_{\nu}
    \;\; = \;\;
        \left[\,0,\infty\,\right)
    \end{equation*}
\item
    Inside a given blackbody radiation cavity at temperature \,$T > 0$,\,
    for each fixed frequency \,$\nu \in (\,0,\infty\,)$,\,
    the probability density function
    induced\footnote{induced by the radiation energy contained inside the given blackbody radiation cavity}
    on the set \,$\mathscr{E}_{\nu}$\, is proportional to the function
    \,$
    [\,0,\infty\,) \xrightarrow{{\color{white}....}} [\,0,\infty\,) :
    {\color{red}
    E \xmapsto{{\color{white}....}} \exp\!\left(\,\overset{{\color{white}.}}{-\,E\,/\,k_{B}\,T}\,\right)
    }
    $,\,
    where \,$k_{B}$\, is the Boltzmann constant.
    In other words, the probability distribution of energy
    inside the given blackbody radiation cavity
    is a \textbf{Boltzmann probability distriubtion}.
    \vskip 0.01cm
    Note in particular that the probability density function here
    is in fact independent of the frequency \,$\nu > 0$,\,
    which of course will imply (as shown below) that the quantity
    \,$\langle\,E(\nu,T)\,\rangle$\,
    will itself be independent of \,$\nu > 0$.\,    
\end{itemize}
\vskip 0.2cm
Under the Reyleigh-Jeans assumptions, we have:
\begin{equation*}
\langle\,E(\nu,T)\,\rangle
\;\; = \;\;
    \dfrac{
        {\color{white}.}
        \displaystyle\int_{0}^{\infty}\;
            E \cdot e^{-E/k_{B}T}
            \;\d E
        {\color{white}.}
        }{
        \overset{{\color{white}.}}{
            \displaystyle\int_{0}^{\infty}\;
                e^{-\,E /k_{B}T}
                \;\d E
            }
        }
\;\; = \;\;
    \cdots
\;\; = \;\;
    \dfrac{
        {\color{white}.}
        \left(\,k_{B}\,T\,\right)^{2}
        {\color{white}.}
        }{
        \overset{{\color{white}.}}{
            k_{B}\,T
            }
        }
\;\; = \;\;
    k_{B}\,T
\end{equation*}
Hence,
\begin{eqnarray*}
\rho(\,\nu,T\,)
& = &
    g(\nu)
    \cdot
    \langle\,E(\nu,T)\,\rangle
\;\; = \;\;
    \dfrac{
        8\,\pi
        }{
        c^{3}
        }
    \cdot
    \nu^{2}
    \cdot
    k_{B}\,T
\\
& = &
    \dfrac{
        8\,\pi\,k_{B}\,T
        }{
        c^{3}
        }
    \cdot
    \nu^{2}
\end{eqnarray*}

          %%%%% ~~~~~~~~~~~~~~~~~~~~ %%%%%

\vskip 0.5cm
\section{Planck's explanation of the universal blackbody spectral energy density via discretization of admissible values of energy}

          %%%%% ~~~~~~~~~~~~~~~~~~~~ %%%%%

\noindent
\textbf{Planck's assumptions}
\begin{itemize}
\item
    Inside a given blackbody radiation cavity,
    for each fixed frequency \,$\nu \in (\,0,\infty\,)$,\,
    the set \,$\mathscr{E}_{\nu}$\, of admissible energies
    among all modes of electromagnetic oscillation
    with frequency \,$\nu$\, is given by:
    \begin{equation*}
    \mathscr{E}_{\nu}
    \;\; = \;\;
        \left\{\;
            \left.
            \overset{{\color{white}1}}{
                n\,h\,\nu \,\in\, \Re
                }
            \;\;\right\vert\;
            n \in \N\cup\{\,0\,\}
            \;\right\},
    \end{equation*}
    where \,$h > 0$\, is the \textbf{\color{red}Planck constant}.
\item
    Inside a given blackbody radiation cavity at temperature \,$T > 0$,\,
    for each fixed frequency \,$\nu \in (\,0,\infty\,)$,\,
    the probability mass function induced
    on the (discrete, countably infinite) set
    \,$\mathscr{E}_{\nu}$\,
    is proportional to the function
    \,$
    \mathscr{E}_{\nu} \xrightarrow{{\color{white}....}} [\,0,\infty) :
    {\color{red}
    n\,h\,\nu
        \xmapsto{{\color{white}....}}
        \exp\!\left(\,\overset{{\color{white}.}}{-\,n\,h\,\nu\,/\,k_{B}\,T}\,\right)
    }
    $,\,
    where \,$k_{B}$\, is the Boltzmann constant.
    \vskip 0.01cm
    So, the probability distribution of energy inside the given blackbody radiation cavity
    assumed by Planck is also a Boltzmann distribution, but it is a Boltzmann distribution
    that is defined on a discrete (countably infinite) set, unlike in the Rayleigh-Jeans case.
\end{itemize}
\vskip 0.2cm
Under Planck's assumptions, we have:
\begin{equation*}
\langle\,E(\nu,T)\,\rangle
\;\; = \;\;
    \dfrac{
        {\color{white}.}
        \overset{\infty}{\underset{n=0}{\sum}}\;
        n\,h\,\nu \cdot e^{-\,n\,h\,\nu/k_{B}T}
        {\color{white}.}
        }{
        \overset{{\color{white}.}}{
            \overset{\infty}{\underset{n=0}{\sum}}\;
            e^{-\,n\,h\,\nu /k_{B}T}
            }
        }
\;\; = \;\;
    \cdots
\;\; = \;\;
    \dfrac{
        h\,\nu
        }{
        {\color{white}.}
        \overset{{\color{white}.}}{
            e^{n\,h\,\nu /k_{B}T} \, - \, 1
            }
        {\color{white}.}
        }
\end{equation*}
Hence,
\begin{eqnarray*}
\rho(\,\nu,T\,)
& = &
    g(\nu)
    \cdot
    \langle\,E(\nu,T)\,\rangle
\;\; = \;\;
    \dfrac{
        8\,\pi
        }{
        c^{3}
        }
    \cdot
    \nu^{2}
    \cdot
    \dfrac{
        h\,\nu
        }{
        {\color{white}.}
        \overset{{\color{white}.}}{
            e^{n\,h\,\nu /k_{B}T} \, - \, 1
            }
        {\color{white}.}
        }
\\
& = &
    \dfrac{
        8\,\pi\,h
        }{
        c^{3}
        }
    \cdot
    \dfrac{
        \nu^{3}
        }{
        {\color{white}.}
        \overset{{\color{white}.}}{
            e^{n\,h\,\nu /k_{B}T} \, - \, 1
            }
        {\color{white}.}
        }
\end{eqnarray*}

          %%%%% ~~~~~~~~~~~~~~~~~~~~ %%%%%

\vskip 0.5cm
\section{Boltzmann's entropy formula \,$S \,=\, k_{B}\log\!\left(\,\Omega\,\right)$}

          %%%%% ~~~~~~~~~~~~~~~~~~~~ %%%%%

Let a thermodynamic system be given, and
let \,$\sigma$\, be a given macroscopic state of the given thermodynamic system.

\vskip 0.4cm
\noindent
We define:
\begin{equation*}
\Omega_{{\color{red}\sigma}}
\;\; := \;\;
    \left(\begin{array}{c}
        \textnormal{number of \textbf{\color{red}microstates}}
        \\
        \textnormal{of the given thermodynamic system}
        \\
        \textnormal{compatible with}
        \\
        \textnormal{the given macroscopic state \,${\color{red}\sigma}$\, of the system}
        \end{array}\right)
\end{equation*}

\vskip 0.4cm
\noindent
We define:
\begin{equation*}
\left(\begin{array}{c}
    \textnormal{the \textbf{\color{red}entropy}}
    \\
    \textnormal{of the given thermodynamic system}
    \\
    \textnormal{at the given macroscopic state \,${\color{red}\sigma}$\, of the system}
    \end{array}\right)
\;\; :=: \;\;
    S_{{\color{red}\sigma}}
\;\; := \;\;
    k_{B}\cdot\log\!\left(\,\Omega_{\sigma}\,\right),
\end{equation*}
where
\,$k_{B} > 0$\,
is the Boltzmann constant.

          %%%%% ~~~~~~~~~~~~~~~~~~~~ %%%%%

\vskip 0.5cm
\section{Thermodynamic definition of absolute temperature}

          %%%%% ~~~~~~~~~~~~~~~~~~~~ %%%%%

Given a thermodynamic system at a given macroscopic state,
the thermodynamic definition of its absolute temperature is:
\begin{equation*}
\dfrac{1}{T}
\;\; := \;\;
    \dfrac{\partial S}{\partial E}\,,
\end{equation*}
where
\,$T$,\, \,$S$,\, and \,$E$\,
are, respectively,
the absolute temperature, the entropy, and the energy
of the given system at the give macroscopic state.
In this sub-section, we explain why this definition
conforms well with our intuitive notion of what temperature is.

\vskip 0.6cm
\begin{center}
\begin{minipage}{4.5in}
% \begin{tcolorbox}[arc=2mm, colback=blue!5!white, colframe=blue!75!black]
% \begin{tcolorbox}[arc=2mm, colback=blue!5!white, colframe=black]
% \begin{tcolorbox}[arc=8px, colback=blue!5!white, colframe=black]
\begin{tcolorbox}[arc=8px, colback=white, colframe=black]
\begin{center}
\vskip 0.1cm
\textbf{Observation \#1:}
\vskip 0.1cm
The absolute temperature \,$T$\,
should be a function of
\,$\dfrac{\partial S}{\partial E}$.\,
\end{center}
\vskip -0.5cm
{\color{white}.}
\end{tcolorbox}
\end{minipage}
\end{center}
\vskip 0.2cm
\noindent
\textbf{Justification of Observation \#1:}
\vskip 0.1cm
\noindent
Our most basic intuitive understanding of temperature
(formalized as the Zeroth Law of Thermodynamics) is that
when you bring two objects into thermal contact,
they will {\color{red}exchange energy} until they reach a state of equilibrium,
at which point they share the same temperature.

\vskip 0.4cm
\noindent
Let us test if our definition satisfies this.
Imagine two isolated systems, $A$ and $B$.
System $A$ has energy $E_{A}$ and entropy $S_{A}(E_{A})$.
System $B$ has energy $E_{B}$ and entropy $S_{B}(E_{B})$.
Now, bring them into thermal contact so they can exchange energy,
but keep the combined system isolated.
The total energy is strictly conserved:
\begin{equation*}
E_{\textnormal{total}}
\;\; = \;\;
    E_{A} \,+\, E_{B}
\;\; = \;\;
    \textnormal{constant}\,,
\end{equation*}
which implies
\begin{equation*}
\d E_{B}
\;\; = \;\;
    -\,\d E_{A}
\end{equation*}
The total entropy of the combined system
is the sum of their individual entropies:
\begin{equation*}
S_{\textnormal{total}}
\;\; = \;\;
    S_{A}(E_{A}) \,+\, S_{B}(E_{B})
\end{equation*}
According to the Second Law of Thermodynamics,
the system will naturally evolve to maximize its total entropy.
{\color{red}At thermal equilibrium}, the total entropy is at a maximum,
hence its first derivative is zero
(${\color{red}\d S_{\textnormal{total}} \,=\, 0}$):
\begin{equation*}
0
\;\; = \;\;
    \d S_{\textnormal{total}}
\;\; = \;\;
    \dfrac{\partial S_{A}}{\partial E_{A}}\,\d E_{A}
    \, + \,
    \dfrac{\partial S_{B}}{\partial E_{B}}\,\d E_{B}
\;\; = \;\;
    \left(\,
        \dfrac{\partial S_{A}}{\partial E_{A}}
        \, - \,
        \dfrac{\partial S_{B}}{\partial E_{B}}
        \right)
    \d E_{A}\,,
\end{equation*}
where the last equality follows by substituting \,$\d E_{B} \,=\, -\,\d E_{A}$.\,
% \begin{equation*}
% dS_{\textnormal{total}}
% \;\; = \;\;
%     \left( \dfrac{\partial S_{A}}{\partial E_{A}} - \dfrac{\partial S_{B}}{\partial E_{B}} \right) dE_{A}
% \;\; = \;\;
%     0
% \end{equation*}
The validity of the above equality
for any {\color{red}arbitrary energy fluctuation} \,$\d E_{A}$\,
implies that the term in the parentheses must be exactly zero,
equivalently
\begin{equation*}
\dfrac{\partial S_{A}}{\partial E_{A}}
\;\; = \;\;
    \dfrac{\partial S_{B}}{\partial E_{B}}
\end{equation*}
The conclusion: At thermal equilibrium,
this specific derivative is exactly the same for both systems.
Because our intuitive definition of equilibrium is ``the state
where $T_A = T_B$'', whatever temperature is,
it must be a direct function of \,$\dfrac{\partial S}{\partial E}$.\,
This completes the justification of Observation \#1.
\qed

\vskip 0.4cm
\noindent
It remains to justify why
the absolute temperature \,$T$\, should be defined to be
the reciprocal of \,$\dfrac{\partial S}{\partial E}$.\,

\vskip 0.6cm
\begin{center}
\begin{minipage}{4.5in}
% \begin{tcolorbox}[arc=2mm, colback=blue!5!white, colframe=blue!75!black]
% \begin{tcolorbox}[arc=2mm, colback=blue!5!white, colframe=black]
% \begin{tcolorbox}[arc=8px, colback=blue!5!white, colframe=black]
\begin{tcolorbox}[arc=8px, colback=white, colframe=black]
\begin{center}
\vskip 0.1cm
\textbf{Observation \#2:}
\end{center}
\vskip -0.1cm
Heat always flows from the system with the smaller value of
\,$\dfrac{\partial S}{\partial E}$\,
to the system with the larger value of
\,$\dfrac{\partial S}{\partial E}$.\,
\vskip -0.2cm
{\color{white}.}
\end{tcolorbox}
\end{minipage}
\end{center}
\vskip 0.2cm
\noindent
\textbf{Justification of Observation \#2:}
\vskip 0.1cm
\noindent
Let's look at the systems before they reach equilibrium.
The Second Law of Thermodynamics dictates that
any spontaneous exchange of energy must strictly increase
the total entropy ($\d S_{\textnormal{total}} > 0$).
Using the same equation as before:
\begin{equation*}
\left(\,
    \dfrac{\partial S_{A}}{\partial E_{A}}
    \, - \,
    \dfrac{\partial S_{B}}{\partial E_{B}}
    \,\right)
\d E_{A}
\;\; = \;\;
    \d S_{\textnormal{total}}
\;\; > \;\;
    0
\end{equation*}
Suppose System $A$ has a larger rate of entropy increase than System $B$, i.e.,
\begin{equation*}
\dfrac{\partial S_{A}}{\partial E_{A}}
\;\; > \;\;
    \dfrac{\partial S_{B}}{\partial E_{B}}
\end{equation*}
For the total entropy change \,$\d S_{\textnormal{total}}$\, to be positive,
\,$\d E_{A}$\, must be positive.
This means energy (heat) must flow into System $A$ and out of System $B$.
Intuitively, if heat is flowing into System $A$,
System $A$ must be the ``colder'' object.
In summary,
\begin{equation*}
\dfrac{\partial S_{A}}{\partial E_{A}}
\; > \;
    \dfrac{\partial S_{B}}{\partial E_{B}}
\quad
\xLeftrightarrow{{\color{white}......}}
\quad
\left\{\begin{array}{c}
    \textnormal{heat flows}
    \\
    \textnormal{from System $B$}
    \\
    \textnormal{into System $A$}
    \end{array}\right\}
\quad
\xLeftrightarrow{{\color{white}......}}
\quad
\left\{\begin{array}{c}
    \textnormal{System $A$}
    \\
    \textnormal{should be colder than}
    \\
    \textnormal{System $B$}
    \end{array}\right\}
\end{equation*}
Therefore, we have established a direct rule:
Heat always flows from the system with the smaller value of
\,$\dfrac{\partial S}{\partial E}$\,
to the system with the larger value of
\,$\dfrac{\partial S}{\partial E}$.\,
This completes the justification of Observation \#2.
\qed

\vskip 0.7cm
\noindent
To align this with our historical, everyday convention that
heat flows from a larger number (higher temperature)
to a smaller number (lower temperature), we must define
temperature \,$T$\, to be a {\color{red}strictly decreasing} function of
\,$\dfrac{\partial S}{\partial E}$.\,
And, of course, the reciprocal is among the simplest strictly decreasing functions.
The foregoing discussion thus leads us to define:
\begin{equation*}
T
\;\; := \;\;
    \left(\,
        \dfrac{\partial S}{\partial E}
        \,\right)^{\!-1}
\end{equation*}

          %%%%% ~~~~~~~~~~~~~~~~~~~~ %%%%%

\vskip 0.5cm
\section{The probability distribution of states of a canonical ensemble is a Boltzmann distribution}

          %%%%% ~~~~~~~~~~~~~~~~~~~~ %%%%%

In this subsection, we establish that the probability distribution of states
for a canonical ensemble is a Boltzmann distribution.

\vskip 0.5cm
\begin{definition}
\mbox{}
\vskip 0.1cm
\noindent
In statistical mechanics, a \textbf{canonical ensemble} is the statistical ensemble
that represents the possible states of a mechanical system
in thermal equilibrium with a heat bath at a fixed temperature.
\end{definition}

\vskip 0.5cm
\begin{proposition}
\mbox{}
\vskip 0.1cm
\noindent
Suppose a canonical ensemble at a fixed temperature $T > 0$\, is given.
Let \,$\Psi$\, denote the micro-state
(considered as a random variable in the sense of probability theory)
of the given canonical ensemble, and
\,$P(\,\Psi = \psi\,)$\,
the probability that the given canonical ensemble
is in the (admissible) micro-state \,$\psi$.\,
Then, we have:
\begin{equation*}
P(\,\Psi = \psi\,)
\;\; = \;\;
    \dfrac{1}{Z}
    \cdot
    \exp\!\left(\,
        -\,
        \dfrac{
            E_{\psi}
            }{
            {\color{white}.}
            k_{B}\,T
            {\color{white}.}
            }
        \,\right),
\quad
\begin{array}{c}
\textnormal{for each admissible micro-state \,$\psi$}
\\
\textnormal{of the given canonical ensemble},
\end{array}
\end{equation*}
where
\,$E_{\psi}$\, is the energy of the micro-state \,$\psi$,\,
\,$k_{B}$\, is the Boltzmann constant, and
\,$Z$\, is a normalization constant.
\end{proposition}
\proof
First, recall that the {\color{red}composite} system
(\,the canonical ensemble together with the heat reservoir \,$R$\,)
is assumed to be an {\color{red}isolated} system.
Hence, the energy of the composite system is conserved:
\begin{equation*}
E_{\psi} \,+\, E_{R}
\;\; \equiv \;\;
    E_{\textnormal{total}}
\;\; = \;\;
    \textnormal{constant}
\end{equation*}
By the {\color{red}Postulate of Equal \textit{A Priori} Probabilities}, we have:
\begin{equation*}
P(\,\Psi = \psi\,)
\;\; = \;\;
    \dfrac{
        \left(\begin{array}{c}
            \textnormal{number of all}
            \\
            \textnormal{accessible micro-states}
            \\
            \textnormal{of the composite system}
            \\
            \textnormal{for which \,$\Psi = \psi$}
            \end{array}\right)
        }{
        \left(\begin{array}{c}
            \textnormal{number of all}
            \\
            \textnormal{accessible micro-states}
            \\
            \textnormal{of the composite system}
            \end{array}\right)
        }
\;\; = \;\;
    \dfrac{
        \Omega_{\Psi}(\psi)
        \times
        \Omega_{R}(E_{R})
        }{
        \Omega_{C}(\,E_{\textnormal{total}}\,)
        }
\;\; = \;\;
    \dfrac{
        \Omega_{R}\!\left(\,E_{\textnormal{total}} \,\overset{{\color{white}.}}{-}\, E_{\psi}\,\right)
        }{
        \Omega_{C}\!\left(\,\overset{{\color{white}.}}{E_{\textnormal{total}}}\,\right)
        }\,,
\end{equation*}
where we have used the following facts:
\,$
E_{\psi} \,+\, E_{R}
\, \equiv \,
    E_{\textnormal{total}}
\, = \,
    \textnormal{constant}
$,\,
and
\begin{equation*}
\left(\begin{array}{c}
    \textnormal{number of all}
    \\
    \textnormal{accessible micro-states}
    \\
    \textnormal{of the composite system}
    \\
    \textnormal{for which \,$\Psi = \psi$}
    \end{array}\right)
\;\; = \;\;
    \underset{\Omega_{\Psi}(\psi)}{\underbrace{\left(\begin{array}{c}
        \textnormal{number of all}
        \\
        \textnormal{accessible micro-states}
        \\
        \textnormal{of the canonical ensemble}
        \\
        \textnormal{for which \,$\Psi = \psi$}
        \end{array}\right)}}
    \times
    \underset{\Omega_{R}(E_{R})}{\underbrace{\left(\begin{array}{c}
        \textnormal{number of all}
        \\
        \textnormal{accessible micro-states}
        \\
        \textnormal{of the heat reservoir}
        \\
        \textnormal{compatible with \,$\Psi = \psi$}
        \end{array}\right)}},
\end{equation*}
and
\,$\Omega_{\Psi}(\psi) \,=\, 1$\,
(since \,$\psi$\, is a micro-state of the canonical ensemble).
\vskip 0.4cm
\noindent
Incidentally, we note that:
\begin{equation*}
\Omega_{C}\!\left(\,E_{\textnormal{total}}\,\right)
\;\; = \;\;
    \underset{\psi}{\sum}\;\,
    \Omega_{\Psi}(\psi)
    \cdot
    \Omega_{R}\!\left(\,E_{\textnormal{total}} \,\overset{{\color{white}.}}{-}\, E_{\psi}\,\right)
\;\; = \;\;
    \underset{\psi}{\sum}\;\,
    \Omega_{R}\!\left(\,E_{\textnormal{total}} \,\overset{{\color{white}.}}{-}\, E_{\psi}\,\right),
\end{equation*}
where the summation is over all accessible micro-states \,$\psi$\,
of the given canonical ensemble.
\vskip 0.4cm
\noindent
By the {\color{red}definition of entropy}, we have:
\begin{equation*}
S_{R}(E_{R})
\;\; = \;\;
    k_{B}
    \cdot
    \log\!\left(\,\overset{{\color{white}.}}{
        \Omega_{R}(E_{R})
        }\,\right)
\;\; = \;\;
    k_{B}
    \cdot
    \log\!\left(\,\overset{{\color{white}.}}{
        \Omega_{R}\!\left(\,E_{\textnormal{total}} \,\overset{{\color{white}1}}{-}\, E_{\psi}\,\right)
        }\,\right)
\end{equation*}
Isolating \,$\Omega_{R}$\, yields:
\begin{equation*}
\Omega_{R}\!\left(\,E_{\textnormal{total}} \,\overset{{\color{white}1}}{-}\, E_{\psi}\,\right)
\;\; = \;\;
    \exp\!\left(\,
        \dfrac{ S_{R}(E_{R}) }{ k_{B} }
        \,\right)
\end{equation*}

\noindent
By Taylor's Theorem, we have:
\begin{eqnarray*}
S_{R}(E_{R})
& = &
    S_{R}\!\left(\,\overset{{\color{white}.}}{
        E_{\textnormal{total}} \,-\, E_{\psi}
        }\,\right)
\\
& \approx &
    S_{R}\!\left(\,\overset{{\color{white}.}}{
        E_{\textnormal{total}}
        }\,\right)
    \, + \,
    \left(
        \left.
        \dfrac{\partial S_{R}}{\partial E_{R}}
        \,\right\vert_{E_{R} \,=\, E_{\textnormal{total}}}
        \right)
    \cdot
    \left(\,\overset{{\color{white}.}}{-\,E_{\psi}}\,\right)
\\
& = &
    S_{R}\!\left(\,\overset{{\color{white}.}}{
        E_{\textnormal{total}}
        }\,\right)
    \, - \,
    \dfrac{E_{\psi}}{T}\,,
\end{eqnarray*}
where the last equality follows from
the thermodynamic definition of absolute temperature
\;$
\dfrac{1}{T}
\, = \,
    \dfrac{
        \partial S
        }{
        \partial E
        }
$.\;

\noindent
Hence,
\begin{eqnarray*}
P(\,\Psi = \psi\,)
& = &
    \cdots
\;\; = \;\;
    \dfrac{
        \Omega_{R}\!\left(\,E_{\textnormal{total}} \,\overset{{\color{white}.}}{-}\, E_{\psi}\,\right)
        }{
        \Omega_{C}\!\left(\,\overset{{\color{white}.}}{E_{\textnormal{total}}}\,\right)
        }
\\
& = &
    \dfrac{1}{\Omega_{C}\!\left(\,\overset{{\color{white}.}}{E_{\textnormal{total}}}\,\right)}
    \cdot
    \exp\!\left(\,
        \dfrac{1}{k_{B}}
        \cdot
        \left(\,
            S_{R}\!\left(\,E_{\textnormal{total}}\,\right)
            \, - \,
            \dfrac{E_{\psi}}{T}
            \,\right)
        \,\right)
\\
& = &
    \dfrac{
        {\color{white}.}
        \exp\!\left(\,\overset{{\color{white}.}}{
            \left.S_{R}\!\left(\,E_{\textnormal{total}}\,\right)\,\right\slash k_{B}
            }\,\right)
        {\color{white}.}
        }{
        \Omega_{C}\!\left(\,\overset{{\color{white}.}}{E_{\textnormal{total}}}\,\right)
        }
    \cdot
    \exp\!\left(\,
        - \,
        \dfrac{
            E_{\psi}
            }{
            {\color{white}.}
            k_{B}\,T
            {\color{white}.}
            }
        \,\right),
\end{eqnarray*}
as required.
\qed

          %%%%% ~~~~~~~~~~~~~~~~~~~~ %%%%%

\vskip 0.5cm
\section{Justification of the use of the Boltzmann probability distriubtion
in the derivation of the Rayleigh-Jeans law}

          %%%%% ~~~~~~~~~~~~~~~~~~~~ %%%%%

We consider the energy contained in the (electromagnetic) radiation
within a blackbody radiation cavity at a certain absolute temperature \,$T$.\,
The electromagnetic radiation obeys Maxwell's equations,
which are linear partial differential equations
(involving partial derivatives of electric and magnetic fields
with respect to space and time).
The linearity of Maxwell's equations implies that their solutions
satisfy the principle of superposition
(i.e., linear combinations of solutions are themselves solutions).
Under appropriate assumptions on cavity geometry and boundary conditions,
we expect that the solution describing the full electromagnetic field
inside the blackbody radiation cavity
can be decomposed into frequency-specific components
(either as a discrete sum or as an integral transform).
In what follows, we will refer to such a frequency-specific component
(which itself can be considered a single-frequency electromagnetic field in its own right)
as a \textbf{mode of oscillation} of the total electromagnetic field.

\vskip 0.3cm
\noindent
Fix a frequency \,$\nu \in (\,0\,,\,\infty\,)$.\,

\begin{equation*}
E_{\textnormal{total}}(\nu)
\;\; = \;\;
    E
    \; + \;
    E_{\textnormal{wall}}
\;\; = \;\;
    \textnormal{constant in time}
\end{equation*}

\begin{equation*}
\Omega_{\textnormal{total}}\!\left(\,\overset{{\color{white}.}}{
    E_{\textnormal{total}}(\nu)
    }\,\right)
\;\; = \;\;
    \underset{E}{\sum}\;\,
    \Omega_{\nu}(E)
    \cdot
    \Omega_{\textnormal{wall}}\!\left(\,\overset{{\color{white}.}}{
        E_{\textnormal{wall}}
        }\,\right)
\;\; = \;\;
    \underset{E}{\sum}\;\,
    \Omega_{\nu}(E)
    \cdot
    \Omega_{\textnormal{wall}}\!\left(\,\overset{{\color{white}.}}{
        E_{\textnormal{total}} \,-\, E
        }\,\right)
\end{equation*}

\vskip 0.3cm
\noindent
Recall the \textbf{Postulate of Equal A Priori Probabilities} from statisticsl mechanics:
\begin{center}
\begin{minipage}{5.5in}
For an isolated macroscopic system in thermodynamic equilibrium
under fixed macroscopic parameters
(such as total energy $E$, volume $V$, and particle number $N$),
all accessible microstates compatible with these constraints
are equally likely to be occupied.
\end{minipage}
\end{center}
\vskip 0.4cm
\noindent
Hence,
\begin{equation*}
P(E)
\;\; \propto \;\;
    \Omega_{\nu}(E)
    \cdot
    \Omega_{\textnormal{wall}}\!\left(\,\overset{{\color{white}.}}{
        E_{\textnormal{total}}-E
        }\,\right)
\end{equation*}

\vskip 0.4cm
\noindent
\textbf{Claim 1}:\;\,
$\Omega_{\nu}(E)$\,
depends only on the frequency \,$\nu$\,
and not on the energy \,$E$.\,
\vskip 0.1cm
\noindent
Proof of Claim 1:\;\;
{\color{red}????}
\vskip 0.1cm
\noindent
This completes the proof of Claim 1.

\vskip 0.4cm
\noindent
By Claim 1, we therefore have:
\begin{equation*}
P(E)
\;\; \propto \;\;
    \Omega_{\textnormal{wall}}\!\left(\,\overset{{\color{white}.}}{
        E_{\textnormal{total}}-E
        }\,\right)
\end{equation*}
Next, by \textbf{\color{red}Boltzmann's entropy formula}
\,${\color{red}
S_{\textnormal{wall}}
\,=\,
    k_{B} \cdot \log\!\left(\,\overset{{\color{white}.}}{
        \Omega_{\textnormal{wall}}
        }\,\right)
}$,\,
we have
\begin{equation*}
\Omega_{\textnormal{wall}}
\;\; = \;\;
    \exp\!\left(\,
        \dfrac{
            S_{\textnormal{wall}}
            }{
            k_{B}
            }
        \,\right)
\end{equation*}
Hence, expressing \,$P(E)$\, in terms of $S_{\textnormal{wall}}$, we have:
\begin{equation*}
P(E)
\;\; \propto \;\;
    \Omega_{\textnormal{wall}}\!\left(\,\overset{{\color{white}.}}{
        E_{\textnormal{total}}-E
        }\,\right)
\;\; = \;\;
    \exp\!\left(\,
        \dfrac{1}{k_{B}}
        \cdot
        S_{\textnormal{wall}}\!\left(\,\overset{{\color{white}.}}{
            E_{\textnormal{total}}-E
            }\,\right)
        \right)
\end{equation*}
By Taylor's Theorem,
\begin{equation*}
S_{\textnormal{wall}}\!\left(\,\overset{{\color{white}.}}{
    E_{\textnormal{total}}-E
    }\,\right)
\;\; \approx \;\;
    S_{\textnormal{wall}}\!\left(\,\overset{{\color{white}.}}{
        E_{\textnormal{total}}
        }\,\right)
    \,-\,
    \left.\dfrac{\partial\,S_{\textnormal{wall}}}{\partial\,E_{\textnormal{wall}}}\right\vert_{E_{\textnormal{total}}}
    \cdot
    E
\end{equation*}
By the \textbf{\color{red}thermodynamic definition of temperature}, we have:
{\color{red}
\begin{equation*}
\dfrac{1}{T_{\textnormal{wall}}}
\;\; := \;\;
    \dfrac{\partial\,S_{\textnormal{wall}}}{\partial\,E_{\textnormal{wall}}}
\end{equation*}
}
Hence,
\begin{eqnarray*}
S_{\textnormal{wall}}\!\left(\,\overset{{\color{white}.}}{
    E_{\textnormal{total}}-E
    }\,\right)
& \approx &
    S_{\textnormal{wall}}\!\left(\,\overset{{\color{white}.}}{
        E_{\textnormal{total}}
        }\,\right)
    \,-\,
    \left.\dfrac{\partial\,S_{\textnormal{wall}}}{\partial\,E_{\textnormal{wall}}}\right\vert_{E_{\textnormal{total}}}
    \cdot
    E
\;\; = \;\;
    S_{\textnormal{wall}}\!\left(\,\overset{{\color{white}.}}{
        E_{\textnormal{total}}
        }\,\right)
    \,-\,
    \dfrac{1}{T_{\textnormal{wall}}} \cdot E
\\
& = &
    S_{\textnormal{wall}}\!\left(\,\overset{{\color{white}.}}{
        E_{\textnormal{total}}
        }\,\right)
    \,-\,
    \dfrac{E}{T}
\end{eqnarray*}
Substituting this into the earlier expression for \,$P(E)$,\, we have:
\begin{eqnarray*}
P(E)
& \propto &
    \Omega_{\textnormal{wall}}\!\left(\,\overset{{\color{white}.}}{
        E_{\textnormal{total}}-E
        }\,\right)
\;\; = \;\;
    \exp\!\left(\,
        \dfrac{1}{k_{B}}
        \cdot
        S_{\textnormal{wall}}\!\left(\,\overset{{\color{white}.}}{
            E_{\textnormal{total}}-E
            }\,\right)
        \right)
\\
& = &
    \exp\!\left(\,
        \dfrac{1}{k_{B}}
        \cdot
        \left(\,
            S_{\textnormal{wall}}\!\left(\,\overset{{\color{white}.}}{
                E_{\textnormal{total}}
                }\,\right)
            \,-\,
            \dfrac{E}{T}
            \,\right)
        \right)
\;\; = \;\;
    \exp\!\left(\,
        \dfrac{
            S_{\textnormal{wall}}\!\left(\,\overset{{\color{white}.}}{
                E_{\textnormal{total}}
                }\,\right)
            }{
            k_{B}
            }
        \right)
    \cdot
    \exp\!\left(\,
        -\,\dfrac{E}{
            {\color{white}.}
            k_{B}\,T
            {\color{white}.}
            }
        \right)
\\
& \propto &
    \exp\!\left(\,
        -\,\dfrac{E}{
            {\color{white}.}
            k_{B}\,T
            {\color{white}.}
            }
        \right)
\end{eqnarray*}

\vskip 0.4cm
\noindent
This completes the proof of the desired proportionality relation:
\begin{equation*}
P(E)
\;\; \propto \;\;
    \exp\!\left(\,
        -\,\dfrac{E}{
            {\color{white}.}
            k_{B}\,T
            {\color{white}.}
            }
        \right)
\end{equation*}

          %%%%% ~~~~~~~~~~~~~~~~~~~~ %%%%%

          %%%%% ~~~~~~~~~~~~~~~~~~~~ %%%%%

          %%%%% ~~~~~~~~~~~~~~~~~~~~ %%%%%

